\chapterwithoutnumber{Introducere}

\section*{Motivația alegerii temei}
\phantomsection
\addcontentsline{toc}{section}{Motivația alegerii temei}

Spațiile verzi și zonele publice deschise sunt printre cele mai vizibile părți ale unui oraș. Ele sunt folosite zilnic pentru deplasare, recreere, sport sau odihnă și influențează direct felul în care este percepută calitatea vieții urbane. Din acest motiv, întreținerea lor nu ține doar de aspectul estetic, ci și de confortul comunității, siguranța utilizării și responsabilitatea față de mediu. În practică, însă, aceste spații sunt afectate frecvent de obiecte abandonate: ambalaje, hârtii, sticle, doze metalice, pungi sau alte resturi de dimensiuni reduse.

La prima vedere, un astfel de obiect poate părea nesemnificativ. Totuși, acumularea deșeurilor în parcuri, pe alei sau în zone publice deschise produce efecte care se observă rapid: crește costul activităților de curățenie, reduce atractivitatea spațiului, creează disconfort pentru cetățeni și poate încuraja repetarea aceluiași comportament. La nivel european, statisticile privind deșeurile municipale sunt urmărite constant deoarece acestea reflectă atât modelele de consum, cât și modul în care colectarea și tratarea deșeurilor sunt organizate administrativ \cite{eurostat2026municipal}. Pentru administratorii spațiilor publice, dificultatea nu constă doar în curățarea locului, ci și în identificarea rapidă a momentelor și zonelor în care apar incidentele.

Monitorizarea clasică se bazează, de regulă, pe patrule periodice, sesizări ale cetățenilor sau camere video folosite mai ales pentru înregistrare pasivă. Aceste metode sunt utile, dar au limite clare. O patrulă nu poate acoperi permanent toate zonele, o sesizare apare de obicei după producerea incidentului, iar o cameră video nu rezolvă problema dacă materialul trebuie verificat manual. O persoană responsabilă poate fi nevoită să urmărească înregistrări lungi pentru a observa un obiect mic, apărut într-o zonă aglomerată sau într-un colț al cadrului. În asemenea situații, oboseala, volumul mare de date și schimbările subtile ale scenei pot duce ușor la omisiuni.

Evoluția viziunii artificiale oferă o soluție potrivită pentru acest tip de problemă. Modelele moderne de detecție a obiectelor pot localiza automat elemente de interes în imagini și cadre video, iar familia YOLO este folosită frecvent în aplicații în care viteza de procesare trebuie echilibrată cu acuratețea \cite{redmon2016yolo}, \cite{ultralytics2023yolov8}. În același timp, seturi de date precum TACO arată că identificarea obiectelor de tip deșeu în contexte reale este o problemă relevantă și studiată în comunitatea de viziune artificială \cite{proenca2020taco}. Diferența importantă, în cazul lucrării de față, este că simpla detectare a unui obiect într-o imagine nu este suficientă. O hârtie poate exista deja în scenă, o pungă poate fi deplasată de vânt, iar o sticlă poate fi confundată parțial cu fundalul. Pentru a vorbi despre un posibil incident, este nevoie de context video: apariția obiectului, prezența unei persoane, evoluția poziției și stabilizarea obiectului în zona monitorizată.

Motivația alegerii temei vine tocmai din această combinație între o problemă practică și un set de tehnologii informatice care pot fi integrate într-un sistem coerent. Lucrarea nu tratează doar recunoașterea unei imagini statice, ci proiectarea unui flux complet: detecția obiectelor de tip deșeu, estimarea materialului, analiza temporală a scenei, salvarea dovezilor locale și validarea incidentelor într-o aplicație web. Prin această abordare, sistemul propus poate reduce timpul necesar verificării manuale și poate ajuta utilizatorul responsabil să se concentreze pe momentele relevante, fără ca decizia finală să fie transferată complet către modelul automat.

\section*{Obiectivele generale ale lucrării}
\phantomsection
\addcontentsline{toc}{section}{Obiectivele generale ale lucrării}

Scopul lucrării este proiectarea, implementarea și evaluarea unui sistem informatic pentru detectarea deșeurilor în spații verzi, utilizând analiza imaginilor video. Sistemul urmărește localizarea obiectelor de tip deșeu, estimarea materialului din care acestea sunt realizate și semnalarea unor posibile incidente prin corelarea detecțiilor cu prezența persoanelor și cu evoluția scenei în timp. Rezultatul practic este o aplicație web locală, cu rol de prototip funcțional, în care utilizatorii pot monitoriza surse video, pot consulta incidentele generate automat și pot confirma, respinge sau arhiva dovezile asociate.

Pentru atingerea acestui scop, sunt urmărite următoarele obiective specifice:

\begin{itemize}
    \item analizarea conceptelor teoretice necesare: viziune artificială, rețele neuronale convoluționale, detecția obiectelor, clasificarea imaginilor și procesarea secvențelor video;
    \item studierea surselor de date disponibile și pregătirea unui set relevant pentru scenarii de monitorizare a deșeurilor în spații verzi;
    \item filtrarea exemplelor necorespunzătoare și organizarea datelor în subseturi distincte de antrenare, validare și testare;
    \item antrenarea prin transfer learning a unui detector de obiecte bazat pe arhitectura YOLOv8s, destinat localizării deșeurilor în imagini;
    \item utilizarea unui clasificator pentru estimarea materialului obiectelor detectate, în categorii precum plastic, hârtie, sticlă, metal și alte materiale;
    \item integrarea unui detector de persoane și proiectarea unei logici temporale care analizează relația dintre prezența umană și apariția unui obiect de tip deșeu;
    \item dezvoltarea unei aplicații web cu autentificare, roluri de acces, monitorizare video, listă de incidente, validare umană și păstrarea locală a dovezilor;
    \item evaluarea componentelor dezvoltate prin metrici specifice, precum precizie, recall, scor F1, mAP@50, mAP@50--95 și acuratețe;
    \item identificarea limitărilor sistemului și formularea unor direcții de îmbunătățire pentru scenarii mai apropiate de utilizarea operațională.
\end{itemize}

\section*{Contribuțiile lucrării}
\phantomsection
\addcontentsline{toc}{section}{Contribuțiile lucrării}

Contribuția principală a lucrării constă în realizarea unui sistem aplicativ care combină viziunea artificială, analiza temporală și dezvoltarea web într-un flux de monitorizare verificabil. Sistemul nu se oprește la afișarea unor casete de detecție peste imagine, ci încearcă să transforme aceste detecții în evenimente gestionabile. În loc să considere fiecare obiect detectat drept incident, aplicația urmărește succesiunea cadrelor și păstrează doar situațiile care merită analizate de utilizatorul responsabil.

Contribuțiile concrete includ pregătirea datelor pentru detecția deșeurilor, antrenarea și integrarea unui detector YOLOv8, folosirea unui clasificator pentru material, integrarea detecției persoanelor, proiectarea unei logici temporale pentru semnalarea incidentelor și dezvoltarea unei aplicații web prin care rezultatele pot fi vizualizate, validate și arhivate. O contribuție importantă este separarea clară dintre sistemul automat și decizia umană. Aplicația nu emite constatări oficiale și nu stabilește vinovății; ea furnizează indicii tehnice și dovezi locale care pot fi analizate de o persoană responsabilă.

Din punct de vedere tehnic, lucrarea reunește componente care, luate separat, sunt relativ cunoscute: modele de învățare automată, procesare video, urmărirea obiectelor în timp, bază de date, autentificare și interfață web. Valoarea proiectului stă în integrarea lor într-un prototip funcțional, orientat către o problemă concretă. Astfel, rezultatul nu este doar un experiment de laborator, ci o aplicație capabilă să gestioneze un flux complet: de la cadru video, la detecție, incident, dovadă și validare.

\section*{Metodologia de cercetare folosită}
\phantomsection
\addcontentsline{toc}{section}{Metodologia de cercetare folosită}

Metodologia lucrării este aplicativ-experimentală. În prima etapă sunt analizate conceptele teoretice necesare pentru înțelegerea soluției: detecția obiectelor, clasificarea imaginilor, învățarea prin transfer, metricile de evaluare, trackingul multi-obiect și analiza temporală a secvențelor video. Fundamentarea teoretică se bazează pe lucrări consacrate din domeniul detecției obiectelor și pe seturi de date utilizate frecvent în viziunea artificială, precum TACO și COCO \cite{proenca2020taco}, \cite{lin2014microsoft}.

În etapa experimentală, datele disponibile sunt analizate și filtrate pentru a obține un set mai apropiat de scopul lucrării. Nu toate imaginile colectate din surse publice sunt utile pentru scenariul urmărit: unele conțin etichete zgomotoase, obiecte prea mici, contexte irelevante sau situații care pot afecta stabilitatea modelului. Din acest motiv, pregătirea datelor este tratată ca parte esențială a proiectului, nu ca o activitate secundară. Setul final este împărțit în subseturi de antrenare, validare și testare, astfel încât performanța să poată fi analizată pe exemple care nu au fost folosite pentru ajustarea modelului.

Antrenarea detectorului se realizează prin transfer learning, nu prin proiectarea unei arhitecturi neuronale de la zero. Această alegere este potrivită pentru o lucrare de licență, deoarece permite valorificarea unor modele preantrenate, capabile să extragă caracteristici vizuale generale. Adaptarea lor la domeniul deșeurilor face posibilă obținerea unui detector specializat cu un volum realist de date și cu resurse hardware accesibile. Clasificarea materialului este tratată separat, fiind aplicată asupra regiunilor de interes generate de detector.

Componenta aplicativă este dezvoltată ca sistem web local, cu posibilitatea de extindere ulterioară către un mediu de server. Backend-ul gestionează autentificarea, baza de date, inferența modelelor, fluxul video și retenția dovezilor, iar frontend-ul oferă interfețe pentru monitorizare, administrare și validarea incidentelor. În cadrul fluxului video, sistemul combină detecția deșeurilor cu detecția persoanelor și cu o logică temporală, deoarece un incident nu poate fi interpretat corect dintr-un singur cadru izolat.

Evaluarea este gândită pe mai multe niveluri. Modelele sunt analizate prin metrici specifice viziunii artificiale, aplicația prin teste funcționale și scenarii de utilizare, iar algoritmul temporal prin clipuri video care permit observarea reacției la apariția obiectelor și la deplasarea persoanelor. Această separare permite identificarea mai clară a punctelor forte și a limitărilor: o eroare poate proveni din detector, din clasificator, din calitatea clipului sau din regulile temporale folosite pentru semnalarea incidentului.

Din punct de vedere al mediului de lucru, experimentele au fost realizate pe un calculator personal cu placă video NVIDIA GeForce RTX 3050 cu 4 GB de memorie, configurație reprezentativă pentru resursele accesibile unui student sau unei instituții mici. Antrenarea și inferența modelelor folosesc biblioteca Ultralytics peste PyTorch, cu accelerare CUDA atunci când placa video este disponibilă și revenire automată pe procesor în caz contrar. Componenta aplicativă folosește FastAPI pentru backend, SQLite pentru persistența locală și OpenCV pentru prelucrarea cadrelor video. Alegerea unor unelte open-source, larg documentate, susține caracterul reproductibil al lucrării: experimentele pot fi reluate pe o configurație similară fără licențe sau servicii externe.

Tot în sprijinul reproductibilității, lucrarea păstrează o separare strictă între artefactele experimentale și cele de producție. Modelele promovate în aplicație sunt stocate în directoare dedicate, distincte de directoarele de experimente, iar seturile de date au împărțiri fixe, cu același seed, astfel încât orice re-rulare a evaluării să producă aceleași subseturi. Pragurile de încredere, dimensiunile de inferență și parametrii algoritmului temporal sunt centralizați în configurația aplicației, ceea ce permite ajustarea lor controlată și documentarea valorilor folosite în teste.

\section*{Limitele lucrării}
\phantomsection
\addcontentsline{toc}{section}{Limitele lucrării}

Lucrarea are câteva limite care trebuie menționate de la început. Performanța sistemului depinde de calitatea imaginilor video: rezoluția camerei, unghiul de filmare, iluminarea, distanța față de obiect, mișcarea persoanelor și nivelul de aglomerație influențează direct rezultatele. Un obiect mic, parțial acoperit sau aflat într-o zonă slab luminată poate fi mai dificil de detectat decât un obiect clar vizibil în prim-plan.

Sistemul nu trebuie interpretat ca o soluție juridică automată. Detectarea unui obiect de tip deșeu și semnalarea unui posibil incident reprezintă indicii tehnice, nu o constatare oficială. Rolul aplicației este de a reduce efortul de observare, de a salva dovezi relevante și de a oferi utilizatorului responsabil un punct de plecare pentru verificare. Decizia finală privind validarea incidentului trebuie să aparțină unei persoane responsabile.

Aria de utilizare este limitată la scenarii apropiate de datele folosite în antrenare și testare. Sistemul este orientat către spații verzi și zone publice deschise, însă nu garantează aceeași performanță în condiții dificile, precum filmare nocturnă, ploaie, zăpadă, aglomerație foarte mare sau camere amplasate la distanțe mari. Extinderea către astfel de situații ar necesita colectarea și etichetarea unor date suplimentare.

O altă limită este caracterul local și experimental al aplicației. Sistemul este gândit pentru rulare locală în contextul lucrării de licență, cu posibilitatea de configurare pentru un deploy ulterior pe server. Pentru utilizare operațională pe termen lung ar fi necesare măsuri suplimentare privind securitatea infrastructurii, administrarea certificatelor, politici detaliate de acces, monitorizare continuă și integrare cu sisteme externe ale autorităților sau administratorilor de spații verzi.

Trebuie menționată și ipoteza de operare a algoritmului temporal: sistemul presupune o cameră fixă, cu un cadru stabil al zonei monitorizate. Această ipoteză este firească pentru scenariul de supraveghere urmărit, însă restrânge tipurile de înregistrări pe care sistemul le poate interpreta corect. Filmările cu camera în mișcare, transfocările sau materialele editate cu tăieturi de scenă rup continuitatea trackingului și fac imposibilă corelarea temporală dintre persoană și obiect. Din acest motiv, atât testarea live, cât și evaluarea pe clipuri din capitolul de testare respectă explicit această ipoteză, iar materialele care o încalcă sunt excluse motivat din scor.

În același registru, distanța de operare influențează direct performanța: un obiect de dimensiuni mici devine, la câțiva metri de cameră, o regiune de doar câteva zeci de pixeli, aflată la limita capacității de detecție. Testele practice efectuate în lucrare arată o funcționare fiabilă la distanțe de până la aproximativ 2--3 metri pentru camera unui telefon; extinderea acestui interval presupune camere cu rezoluție mai mare sau decuparea inteligentă a zonelor de interes, direcții discutate la finalul lucrării.

\section*{Structura lucrării}
\phantomsection
\addcontentsline{toc}{section}{Structura lucrării}

Lucrarea este organizată în cinci capitole. Primul capitol prezintă problema deșeurilor în spații publice, necesitatea monitorizării automate și termenii utilizați în continuare. Capitolul al doilea descrie fundamentele teoretice și tehnologiile folosite: viziunea artificială, detecția obiectelor, clasificarea imaginilor, trackingul, analiza temporală și tehnologiile web. Capitolul al treilea este dedicat datelor, antrenării modelelor și metodologiei experimentale. Capitolul al patrulea prezintă proiectarea și implementarea sistemului, incluzând arhitectura aplicației, integrarea pipeline-ului ML, algoritmul temporal, backend-ul, baza de date și interfața web. Capitolul al cincilea tratează testarea, rezultatele obținute, analiza erorilor și direcțiile de îmbunătățire. Lucrarea se încheie cu sinteza concluziilor, contribuțiile personale și limitările rămase.

Ordinea capitolelor urmează parcursul firesc al proiectului: de la delimitarea problemei și a conceptelor, prin pregătirea datelor și antrenarea modelelor, către integrarea lor într-o aplicație completă și, în final, către verificarea experimentală a întregului sistem. Fiecare capitol se sprijină pe cele anterioare, astfel încât cititorul să poată urmări atât raționamentul tehnic, cât și deciziile practice care au condus la forma finală a aplicației.

Lucrarea poate fi parcursă, de altfel, pe două planuri complementare. Primul este planul de învățare automată: datele, antrenarea modelelor și metricile de evaluare, tratate în principal în capitolele 3 și 5. Al doilea este planul ingineriei aplicației: arhitectura, fluxurile de utilizare, securitatea și gestionarea dovezilor, tratate în capitolul 4. Capitolele 1 și 2 leagă cele două planuri prin cerințele domeniului și prin fundamentele teoretice comune. Această dublă perspectivă reflectă natura proiectului: un model performant fără o aplicație în jurul lui rămâne un experiment de laborator, iar o aplicație fără un model evaluat riguros rămâne o simplă demonstrație --- lucrarea urmărește să le susțină pe amândouă, cu aceeași atenție pentru reproducibilitate.

Pe lângă evaluarea finală, dezvoltarea aplicației a fost însoțită permanent de verificări automate și manuale. Proiectul include o suită de teste automate pentru componentele backend-ului --- autentificare, reguli de acces, operații pe incidente --- rulată după fiecare modificare importantă, precum și scripturi de verificare rapidă pentru pipeline-ul de inferență și pentru fluxul video complet. Acest mod de lucru a permis identificarea timpurie a unor defecte subtile, cum ar fi pierderea clipului de dovadă atunci când două incidente apăreau la interval foarte scurt, și corectarea lor înainte de evaluarea finală. Într-un sistem cu multe componente interdependente, în care o modificare a algoritmului temporal poate afecta indirect salvarea dovezilor sau raportarea statisticilor, această disciplină de verificare s-a dovedit la fel de importantă ca alegerea modelelor în sine. Capitolul de testare reia aceste aspecte în mod sistematic, de la metricile modelelor până la validarea funcțională a aplicației web în ansamblu.
